% © 2026 Ing. David Jaroš — CC BY-NC-ND 4.0
%
% This work is licensed under a Creative Commons Attribution-NonCommercial-NoDerivatives
% 4.0 International License (CC BY-NC-ND 4.0).
%
% License History: Earlier drafts (up to v0.3) were released under CC BY 4.0.
% From v0.4 onward, all material is released under CC BY-NC-ND 4.0 to protect
% the integrity of the theoretical work during ongoing academic development.
%
% See LICENSE.md for full license text.
%
% Track: quantum_ubt — Cosmological solutions
% File: research_tracks/quantum_ubt/frw_from_ubt.tex
% Purpose: First step toward deriving FRW cosmology from the UBT Θ field.
%          Investigates the FRW ansatz Θ(t) = a(t)Θ₀ and checks whether
%          Friedmann equations follow from the UBT metric chain.
% Status: [Prop.] — ansatz formulated; Friedmann equations follow at [Prop.] level.
%         Full derivation from S[Θ] is [OPEN] (GAP-C).
% Date: 2026-05-17

\documentclass[12pt,a4paper]{article}
\usepackage[a4paper,margin=2.5cm]{geometry}
\usepackage{amsmath,amssymb,amsthm,booktabs}
\usepackage{tcolorbox}
\tcbuselibrary{skins,breakable}
\usepackage{hyperref}

\newtcolorbox{statusbox}[1][]{%
  colback=gray!8!white,colframe=black!60,arc=3pt,boxrule=1pt,
  fonttitle=\bfseries,title=Status Summary,#1}
\newtcolorbox{propbox}[1][]{%
  colback=green!4!white,colframe=green!55!black,arc=3pt,boxrule=1pt,
  fonttitle=\bfseries,title=Proposition (to be verified),#1}
\newtcolorbox{openbox}[1][]{%
  colback=yellow!5!white,colframe=orange!55!black,arc=3pt,boxrule=1pt,
  fonttitle=\bfseries,title=Open Problem,#1}

\newtheorem{theorem}{Theorem}[section]
\newtheorem{lemma}[theorem]{Lemma}
\newtheorem{proposition}[theorem]{Proposition}
\newtheorem{corollary}[theorem]{Corollary}

\title{FRW Cosmology from UBT:\\[4pt]
       First Steps — The Homogeneous Ansatz}
\author{Ing.\ David Jaro\v{s}}
\date{2026-05-17}

\begin{document}
\maketitle

\begin{abstract}
We investigate whether the UBT framework reproduces Friedmann--Robertson--Walker
(FRW) cosmology from the $\Theta$ field.  A homogeneous isotropic ansatz
$\Theta(t) = a(t)\,\Theta_0$ (no spatial dependence) is proposed and checked
against the UBT metric chain
$g_{\mu\nu} = \mathrm{Re}[\mathrm{Tr}(\partial_\mu\Theta\cdot\partial_\nu\Theta^\dagger)]/\mathcal{N}$.
Preliminary analysis shows that this ansatz produces a conformally flat metric
proportional to $a(t)^2\,\eta_{\mu\nu}$; recovery of the full FRW metric requires
an off-diagonal structure in $\partial_\mu\Theta$.  A refined ansatz is proposed
and Friedmann equations are argued to follow at \textbf{[Prop.]} level.
The de Sitter limit is checked for consistency with the Casimir energy.
GAP-C status: \textbf{OPEN} with active first steps.
\end{abstract}

\tableofcontents

%──────────────────────────────────────────────────────────────────────────────
\section{UBT Metric Chain}
\label{sec:metric_chain}
%──────────────────────────────────────────────────────────────────────────────

From \texttt{canonical/gr\_closure/step1\_metric\_bridge.tex}~[L1]:
\begin{equation}
  \label{eq:metric_chain}
  g_{\mu\nu} = \frac{1}{\mathcal{N}}\,
  \mathrm{Re}\bigl[\mathrm{Tr}(\partial_\mu\Theta\cdot\partial_\nu\Theta^\dagger)\bigr].
\end{equation}
The normalisation $\mathcal{N}$ is fixed by the large-distance flat limit
$g_{\mu\nu} \to \eta_{\mu\nu}$.

The Schwarzschild metric is recovered from a spherically symmetric ansatz
$\Theta_0(r,t)$ (see \texttt{papers/UBT\_GR\_Submission.tex}~[L1]).
GAP-C (\texttt{WHAT\_IS\_PROVED.md}) asks whether an FRW ansatz exists.

%──────────────────────────────────────────────────────────────────────────────
\section{Homogeneous Isotropic Ansatz}
\label{sec:frw_ansatz}
%──────────────────────────────────────────────────────────────────────────────

\subsection{First ansatz: $\Theta(t) = a(t)\,\Theta_0$}

For a homogeneous isotropic universe, we try
\begin{equation}
  \label{eq:ansatz1}
  \Theta(t) = a(t)\,\Theta_0,
\end{equation}
where $\Theta_0\in\mathbb{B}$ is a constant biquaternion and
$a(t)$ is the scale factor.

With no spatial dependence, $\partial_i\Theta = 0$ for $i=1,2,3$, so:
\begin{align}
  \partial_t\Theta &= \dot a(t)\,\Theta_0, \\
  \partial_i\Theta &= 0.
\end{align}

Substituting into~(\ref{eq:metric_chain}):
\begin{align}
  g_{tt} &= \frac{\dot a^2}{\mathcal{N}}
             \mathrm{Re}[\mathrm{Tr}(\Theta_0\cdot\Theta_0^\dagger)]
           = \frac{\dot a^2 \|\Theta_0\|^2}{\mathcal{N}}, \\
  g_{ti} &= 0, \\
  g_{ij} &= 0.
\end{align}

This gives a metric $g_{\mu\nu} = \mathrm{diag}(\dot a^2 C, 0, 0, 0)$
which is \emph{degenerate} in the spatial directions.
Ansatz~(\ref{eq:ansatz1}) does \emph{not} produce the FRW metric.

\subsection{Refined ansatz: $\Theta(x,t) = a(t)\,e^{ix^k e_k}\,\Theta_0$}

The spatial metric components require spatial gradients.
We consider the refined ansatz
\begin{equation}
  \label{eq:ansatz2}
  \Theta(x,t) = a(t)\,\exp\!\bigl(i\,x^k e_k\bigr)\,\Theta_0,
\end{equation}
where $e_k = \sigma_k/2$ are the $\mathrm{SU}(2)$ generators
and $x^k = (x^1, x^2, x^3)$.
This is a coherent superposition of three winding modes in the spatial directions.

At leading order in small $|x|$:
\begin{align}
  \partial_t\Theta &= \dot a\,e^{ix^k e_k}\,\Theta_0, \\
  \partial_j\Theta &= a(t)\,i\,e_j\,e^{ix^k e_k}\,\Theta_0
  \approx i\,a(t)\,e_j\,\Theta_0 + \mathcal{O}(x).
\end{align}

The spatial metric at $x = 0$:
\begin{align}
  g_{ij}\big|_{x=0}
  &= \frac{a^2}{\mathcal{N}}\,
  \mathrm{Re}\bigl[\mathrm{Tr}(e_i\,\Theta_0\cdot\Theta_0^\dagger e_j^\dagger)\bigr].
\end{align}
For $\Theta_0 = \mathbf{1}_2$ and $e_k = \sigma_k/2$,
$\mathrm{Tr}(e_i e_j^\dagger) = \frac{1}{2}\delta_{ij}$, giving
\begin{equation}
  g_{ij}\big|_{x=0} = a^2(t)\,\delta_{ij}\,C', \qquad
  C' = \frac{1}{2\mathcal{N}}.
\end{equation}
Choosing $\mathcal{N}$ such that $C'=1$ at $a=1$:

\begin{propbox}
The refined ansatz~(\ref{eq:ansatz2}) with $\Theta_0 = \mathbf{1}_2$ gives,
at linear order in $x$, the spatially flat FRW metric
\[
  ds^2 = -f(\dot a)^2\,dt^2 + a^2(t)\,(dx^2 + dy^2 + dz^2),
\]
where $f(\dot a)$ is fixed by the $tt$ component.
Full verification away from $x=0$ and the Friedmann equations are~\textbf{[Prop.]}
pending detailed computation.
\end{propbox}

%──────────────────────────────────────────────────────────────────────────────
\section{Friedmann Equations from UBT}
\label{sec:friedmann}
%──────────────────────────────────────────────────────────────────────────────

\begin{proposition}[Friedmann equations, \textbf{[Prop.]}]
\label{prop:friedmann}
If the metric~(\ref{eq:metric_chain}) with ansatz~(\ref{eq:ansatz2})
reproduces $g_{\mu\nu} = \mathrm{diag}(-1, a^2, a^2, a^2)$ globally,
then the UBT field equation $\nabla^\dagger\nabla\Theta = \kappa\mathcal{T}$
reduces, in the homogeneous isotropic sector, to the Friedmann equations
\begin{align}
  \label{eq:friedmann1}
  H^2 &= \frac{8\pi G}{3}\,\rho, \\
  \label{eq:friedmann2}
  \dot H &= -4\pi G\,(\rho + p),
\end{align}
where $H = \dot a/a$ is the Hubble parameter.
\end{proposition}

\textit{Argument}: Equations~(\ref{eq:friedmann1})--(\ref{eq:friedmann2})
follow directly from the Einstein equations $G_{\mu\nu} = 8\pi G\,T_{\mu\nu}$
(proved at [L1] in \texttt{canonical/gr\_closure/}) with the FRW
metric, by standard GR computation.  The non-trivial step is showing that
the UBT ansatz~(\ref{eq:ansatz2}) produces the FRW metric globally, which
remains at \textbf{[Prop.]} status.

%──────────────────────────────────────────────────────────────────────────────
\section{de Sitter Limit}
\label{sec:de_sitter}
%──────────────────────────────────────────────────────────────────────────────

\subsection{Exponential expansion ansatz}

For $a(t) = e^{Ht}$ (de Sitter expansion), the Friedmann equations give:
\[
  H^2 = \frac{\Lambda}{3}, \qquad \rho + p = 0.
\]

In UBT the vacuum energy is sourced by the Casimir contribution of the
KK sector on $S^1_\psi$.  From \texttt{tools/casimir\_B\_check.py}
(see also \texttt{research\_tracks/T3\_ALPHA/mellin\_insertion\_B.tex}):
\[
  E_{\rm Casimir}(T^3, R_\psi = 1) \approx -0.0617\quad (\mathrm{in\,units\,of}\,R_\psi^{-3}).
\]

\subsection{Consistency with cosmological constant}

The de Sitter Hubble parameter would be
\[
  H_{\rm UBT}^2 \sim \frac{8\pi G}{3}\,|E_{\rm Casimir}|\,R_\psi^{-3}
  \sim \frac{0.517}{\ell_{\rm Pl}^2\cdot 3} \sim M_{\rm Pl}^2.
\]
This is of Planck-scale magnitude, vastly larger than the observed
$\Lambda_{\rm obs} \sim 10^{-122}\,M_{\rm Pl}^4/M_{\rm Pl}^4$.
This is the standard cosmological constant problem, not specific to UBT.

\begin{openbox}
The de Sitter limit of UBT and the smallness of $\Lambda_{\rm obs}$
require a cancellation mechanism.  This is~\textbf{[OPEN]} and tracked
as GAP-C in \texttt{WHAT\_IS\_PROVED.md}.
\end{openbox}

%──────────────────────────────────────────────────────────────────────────────
\begin{statusbox}
FRW ansatz $\Theta = a(t)\Theta_0$ (simple): does not produce FRW metric [checked]\\
Refined ansatz $\Theta = a(t)e^{ix^k e_k}\Theta_0$: FRW metric at $x=0$ [Prop.]\\
Global FRW metric from refined ansatz: \textbf{[Prop.]} — verification needed\\
Friedmann equations from UBT: \textbf{[Prop.]} — conditional on metric [Prop.]\\
de Sitter limit: $H^2 \sim M_{\rm Pl}^2$ (Planck-scale, as expected) [checked]\\
Cosmological constant problem: \textbf{[OPEN]} (standard issue)\\
GAP-C status: \textbf{OPEN} — active first steps; route identified
\end{statusbox}
%──────────────────────────────────────────────────────────────────────────────

\end{document}
